\subsection{A Shared Profile for Training and Execution}
\label{sec:impl:profile_revision}
The revised selector uses one feature function for dataset construction and
execution. Its eight inputs are the four profile scores, node count, triple
count, directed simple-projection density, and the number of detected
violations. Labels describe injected defects; they do not enter this feature
function. The document identifier and allowed relations are public task
metadata. The checkpoint and scaler carry a feature-version identifier;
incompatible weights cause an explicit uniform-prior fallback.

The source-field check recognizes a nonempty value after an explicit allowed
field label at a document, line, or sentence boundary. A missing field is
reported only when such a value is present in the source. Conflicting values
under a repeated label cause abstention. A relation vocabulary alone does not
make every relation mandatory. These source-field reports enter the violation
count and candidate pool. They do not change the meaning of node incidence:
a partial star graph can still have incidence score 100.

Hierarchy checks distinguish downward relations (management and supervision)
from upward membership. Ambiguous relation meanings are left unchecked.
Endpoint membership is checked against declared entities only for endpoints
explicitly typed as entity references; literal field values are not dangling
entities. The empty-graph guard applies to both individual and bundled edits.
A supplied candidate bypasses the low-probability and no-violation early stops,
so each proposed change receives a trial evaluation.

For profile quality $Q\in[0,100]$, hard-violation count $h$, operation cost
$c(a)$, scope $s(a)$, and proposal confidence $z(a)$, the default utility is
\begin{equation}
\begin{split}
U(a)={}&\frac{Q(G_a)-Q(G)}{100}
+0.20\frac{\max(0,h(G)-h(G_a))}{\max(1,h(G))}\\
&-0.35c(a)+0.05\log\!\left(\max(\pi_{s(a)},0.05)+10^{-6}\right)
+0.02z(a).
\end{split}
\end{equation}
Here $G_a$ is the trial graph. Delete, retype, and add cost 0.30, 0.16,
and 0.06; bundle costs sum their operations. A feasible edit is accepted only
if its utility is positive and maximal among the evaluated edits. The
four quality terms retain equal weights. Feasibility preserves each satisfied
lower bound, does not worsen a deficient dimension, limits a checked
component's regression to two points, and enforces the density bound.
These conditions are local constraints, rather than a monotonicity result for
reference F1. Pairwise redundancy assessment costs $O(m^2)$ for $m$ triples;
$T$ rounds with at most $K$ trial candidates cost $O(TKm^2)$, apart from
source-string operations. This differs from the main filter's linear
candidate scan.

We predefine a second utility by subtracting
$0.05\log(1/3+10^{-6})$ from the default expression. This makes the prior
term zero under a uniform distribution. It leaves the within-step ordering
unchanged for a fixed router, but changes which candidates have positive
utility. Neither formula is tuned on the archived proposals.
